\chapter{Transurethral Resection of the Prostate (TURP)}
\section*{What Is This Test or Procedure?}
TURP removes obstructing prostate tissue through a resectoscope passed through the urethra.
\section*{Why This Test or Procedure Is Ordered}
It treats bothersome or complicated urinary obstruction from benign prostatic enlargement when medicines are inadequate or complications develop.
\section*{How This Test or Procedure Is Performed}
Under regional or general anesthesia, an electrically powered or laser instrument removes prostate tissue in small pieces. A urinary catheter remains temporarily.
\section*{Risks and Recovery}
Bleeding, infection, urinary retention, urethral stricture, incontinence, erectile or ejaculatory changes, and electrolyte disturbance are possible. Urinary symptoms improve gradually during healing.
\section*{References}
\begin{enumerate}\item MedlinePlus. Transurethral Resection of the Prostate. U.S. National Library of Medicine; 2025.\item Current Medical Diagnosis and Treatment. McGraw-Hill; 2025.\end{enumerate}
\clearpage
